\section{Implementation}
\label{sec:system}

BeyondNN~0.1.0 is a Python library for PyTorch released under the Apache-2.0 licence (\texttt{pip install beyondnn}; source at \url{https://github.com/NikolasRoufas/beyondnn}). It has 75 modules, is tested on Python 3.10, 3.12 and 3.14 with PyTorch~$\ge$2.3, and \texttt{import beyondnn} does not import PyTorch. We describe only what the evaluation uses; all names are public API.

\paragraph*{Layers}
\texttt{core} handles tracing, hooks, sites, sample identity and persistence; \texttt{schema} and \texttt{provenance} hold the immutable, versioned records, statuses, claims and model fingerprints. Five modules produce or assess evidence: \texttt{attribution}, \texttt{interventions}, \texttt{faithfulness}, \texttt{concepts} and \texttt{audits}. \texttt{explain} composes the structured WHY. Every layer below \texttt{audits} produces records; the audit only reads them. Each record kind is versioned and older versions are migrated on load.

\paragraph*{Instrumentation}
\texttt{bnn.trace(model, x, sites=...)} runs one forward pass and returns a trace; hooks exist only during a recording, and the user's model is never modified. Tracing refuses situations in which a record could be silently wrong (foreign forward hooks, models spread over several devices, structural changes during recording).

\paragraph*{Evidence producers}
\texttt{bnn.attribute} records \textsc{attributed} evidence (native gradient, input$\times$gradient and IG; Captum methods through an optional adapter, verified with Captum~0.9). \texttt{bnn.intervene} records \textsc{interventional} effects of the declared operations \texttt{zero}, \texttt{constant}, \texttt{patch} and \texttt{direction} on sites, units or feature directions. The \texttt{faithfulness} module implements selection-based tests: \texttt{F.top\_k} turns an attribution into a selection (descending absolute score, ties broken by index, optionally among declared eligible units), \texttt{F.comprehensiveness} and \texttt{F.sufficiency} declare the test of~\eqref{eq:test}, \texttt{F.replacement} declares a named replacement tensor (for example a training mean, a resampled activation or a \texttt{[MASK]} embedding), and \texttt{F.controls} declares matched random controls with a recorded seed. There is no implicit zero replacement.

\paragraph*{Plans, audits and persistence}
An audit plan is itself a content-addressed record. It declares the checkpoint, the samples, the claims, each claim's requirement (assessment policy, whether controls are mandatory, alternative threshold criteria), the invariances the claim asserts, the configuration roles and counterexample caps. \texttt{bnn.audit(evidence, plan=plan)} returns a deterministic report with sorted keys that embeds the plan, a digest of the evidence identifiers and its producer (BeyondNN version and audit-semantics version). \texttt{AU.save\_evidence} and \texttt{AU.load\_evidence} save and reload exactly the traces an audit ingests, with full re-validation; \texttt{AU.verify\_report(report, evidence, plan)} re-derives a stored report and raises an error on any difference. Concept evidence follows the same pattern: an encoding test against random directions and label permutations, and a use test that removes the concept direction and compares the target change with a declared minimum in the target's own scale.
